\begin{abstract}
The transition from IPv4 to IPv6 changes how residential Internet of Things (IoT) devices are exposed. Under IPv4, NAT-based residential deployments commonly discarded unsolicited inbound traffic by default, giving devices an implicit perimeter and limiting their direct reachability from the public Internet. IPv6 removes this form of implicit protection. Devices acquire globally routable addresses and, depending on the router configuration, may become directly reachable from the public Internet without any change to their own security properties.

Prior work has studied this shift only in fragments, including residential firewall behavior, IPv6 host discovery, device identifiability, and firmware weaknesses. The path from a globally addressed device to externally accessible services has received comparatively limited attention under controlled residential conditions. This thesis addresses that gap by deploying an ESP32 device running custom firmware on the lwIP stack and exposing deliberately weak HTTP and Telnet services on a real residential IPv6 connection. The device is observed across five experimental scenarios.

The results show that exposure was governed primarily by the residential firewall rather than by changes in the device itself. With all device-side conditions held constant, disabling the perimeter moved the same host from externally inaccessible services to direct external access to its HTTP endpoint, hidden administrative interface, and Telnet service. Filtering, however, reduced accessibility but not observability. The device remained reachable through ICMPv6 and identifiable through its EUI-64-derived address even when all inbound TCP traffic was blocked. In contrast, the IPv4/NAT baseline kept the same services unreachable without port forwarding. These findings suggest that IPv6 does not inherently increase device insecurity, but it does alter the conditions under which vulnerable services become reachable. In the scenarios examined, exposure depended largely on network-level configuration choices rather than on actions taken by end users.
\end{abstract}
